\section{Artin--Schreier--Witt theory}
\label{appendix:asw_development}

We retain the notation for Witt vectors from the beginning of this appendix.
For an $\F_p$-algebra $A$, Frobenius acts componentwise on $W_r(A)$.

\subsection{The ASW operator and bookkeeping}

\begin{definition}\label{definition:asw_operator}
The \textbf{Artin--Schreier--Witt operator} is
\[
\wp:=F-\mathrm{id},\qquad
\wp\vec u=F\vec u\ominus\vec u.
\]
It is an endomorphism of the additive group of $W_r(A)$ and commutes with
truncation and Verschiebung.
\end{definition}

\begin{lemma}[Witt bookkeeping]\label{lemma:witt_bookkeeping}
Let $A$ be an $\F_p$-algebra.
\begin{enumerate}[label=(\roman*)]
    \item The truncation $t:W_r(A)\to W_{r-1}(A)$ is a ring homomorphism
    commuting with $F$, hence with $\wp$, and
    $\ker(t)=V^{r-1}A$.
    \item The map $V^{r-1}:A\to W_r(A)$ is additive and
    $\wp\circ V^{r-1}=V^{r-1}\circ\wp$.
    \item For every $\vec a\in W_r(A)$ and $c\in A$,
    \[
    \vec a\oplus V^{r-1}(c)
    =(a_0,\dots,a_{r-2},a_{r-1}+c).
    \]
\end{enumerate}
\end{lemma}

\begin{proof}
The first two assertions follow from
\Cref{prop:witt_multiplication_by_p}. For (iii), consider the two operations
\[
\begin{aligned}
 \Phi(\vec X,C)
 &:=\vec X\oplus V^{r-1}(C)
   =\bigl(S_n(\vec X;0,\dots,0,C)\bigr)_{0\leq n\leq r-1},\\
 \Psi(\vec X,C)
 &:=(X_0,\dots,X_{r-2},X_{r-1}+C).
\end{aligned}
\]
Both are given by universal polynomials with integer coefficients, as required
in \Cref{lemma:ghost_principle}. Over
$\Z[X_0,\dots,X_{r-1},C]$, for $n<r-1$ we have
\[
 w_n\bigl(\Phi(\vec X,C)\bigr)
 =w_n(\vec X)+w_n\bigl(V^{r-1}(C)\bigr)
 =w_n(\vec X)
 =w_n\bigl(\Psi(\vec X,C)\bigr).
\]
For the final ghost component,
\[
\begin{aligned}
 w_{r-1}\bigl(\Phi(\vec X,C)\bigr)
 &=w_{r-1}(\vec X)+p^{r-1}C\\
 &=\sum_{i=0}^{r-2}p^iX_i^{p^{r-1-i}}
   +p^{r-1}(X_{r-1}+C)\\
 &=w_{r-1}\bigl(\Psi(\vec X,C)\bigr).
\end{aligned}
\]
Thus $w\circ\Phi=w\circ\Psi$, so
\Cref{lemma:ghost_principle} gives $\Phi=\Psi$ over every commutative ring,
and in particular over $A$.
\end{proof}

The last assertion is special to the top Verschiebung slice. Witt addition is
not componentwise in lower positions, and the carry polynomials in
\Cref{eq:witt_addition_shape} cannot in general be discarded.

\subsection{Torsors and classification}
The Witt addition polynomials make $\A^r_{\F_p}$ into the additive Witt
vector group scheme $\bW_r$. We distinguish this group scheme from its
groups of points: for an $\F_p$-algebra $A$, its group of $A$-valued points
is $\bW_r(A)=W_r(A)$. If $X$ is an $\F_p$-scheme, we write
$W_r(\F_p)_X$ for the constant group scheme over $X$ associated with the
abstract group $W_r(\F_p)$. The operator $\wp=F-\mathrm{id}$ is an
endomorphism of $\bW_r$ and fits into an exact sequence of group schemes for
the \'etale topology
\begin{equation}\label{eq:asw_lang_sequence}
0\longrightarrow W_r(\F_p)_{\F_p}\cong(\Z/p^r\Z)_{\F_p}
\longrightarrow\bW_r\xrightarrow{\wp}\bW_r\longrightarrow0.
\end{equation}

The following is the Witt-vector analogue of
\Cref{lemma:classical_cyclic_equations_etale_torsors}. Its coordinate proof
requires the triangular form of Witt addition because, unlike the two
classical equations, Witt subtraction is not componentwise.

\begin{lemma}[ASW covers are \'etale torsors]
\label{lemma:asw_etale_torsor}
Let $A$ be an $\F_p$-algebra and $\vec g\in W_r(A)$. Let
$\vec X=(X_0,\dots,X_{r-1})$ denote the universal Witt vector of length $r$
over $A$, so that $A[\vec X]=A[X_0,\dots,X_{r-1}]$. The fiber of $\wp$
over $\vec g$ is
\[
T_{\vec g}=\Spec B,
\qquad
B:=A[\vec X]\big/\bigl(\wp\vec X\ominus\vec g\bigr),
\]
where the parenthesized Witt vector denotes the ideal generated by its $r$
coordinate polynomials.
Here $\vec g\in W_r(A)=\bW_r(A)$ is viewed as a morphism
$\Spec A\to\bW_r$, and $T_{\vec g}$ is defined by the Cartesian square
\[
\begin{tikzcd}
T_{\vec g} \arrow[r] \arrow[d] & \bW_r \arrow[d,"\wp"] \\
\Spec A \arrow[r,"\vec g"'] \arrow[dr] & \bW_r \arrow[d] \\
& \Spec\F_p,
\end{tikzcd}
\]
where the maps to $\Spec\F_p$ are the structure maps.
The algebra $B$ is finite free over $A$, with basis
$\{\prod_nX_n^{a_n}:0\leq a_n\leq p-1\}$, and is \'etale over $A$.
Translation by $W_r(\F_p)\cong\Z/p^r\Z$, explicitly
$\vec X\mapsto\vec X\oplus\vec c$, makes $T_{\vec g}$ a
$\Z/p^r\Z$-torsor over $\Spec A$.
\end{lemma}

\begin{proof}
By the triangular shape of Witt addition
(\Cref{eq:witt_addition_shape}), the $n$-th equation in
$\wp\vec X\ominus\vec g=0$ has the form
\[
X_n^p-X_n-g_n+C_n'(X_0,\dots,X_{n-1};g_0,\dots,g_{n-1})=0.
\]
Starting with $A$, first adjoin $X_0$ subject to the equation of index $0$.
Then, for each $1\leq n<r$, adjoin $X_n$ subject to the displayed equation
of index $n$; its coefficients lie in the algebra obtained at the preceding
step. Because this equation is monic of degree $p$ in $X_n$, the $n$-th step
is free of rank $p$, with basis $1,X_n,\dots,X_n^{p-1}$. After all $r$ steps
we obtain $B$, finite free over $A$ with the stated basis. The Jacobian matrix
is lower triangular with diagonal entries $-1$, so $B/A$ is \'etale.
Moreover, this cover is the
pullback of $\wp:\bW_r\to\bW_r$ along the section $\vec g$. The kernel of
the group-scheme morphism $\wp$ is the constant group scheme
$W_r(\F_p)_{\F_p}$, characterized by $F\vec u=\vec u$. Translation by the
base change of this kernel to $\Spec A$ preserves the fiber because, for
$\vec c\in\ker(\wp)$,
\[
\wp(\vec X\oplus\vec c)=\wp(\vec X)\oplus\wp(\vec c)=\wp(\vec X).
\]
Any two points in the same fiber differ by a unique such $\vec c$, so this
translation action makes the pullback a torsor.
\end{proof}

\begin{theorem}[Artin--Schreier--Witt theory;
{\cite[\S3--4]{Thomas2005}}]
\label{theorem:asw_theory}
Let $M$ be a field of characteristic $p$, and let $r\geq1$. For
$\vec g\in W_r(M)$, let
\[
T_{\vec g}:=
\Spec\!\left(M[\vec X]\big/\bigl(\wp\vec X\ominus\vec g\bigr)\right)
\]
be the $\Z/p^r\Z$-torsor of \Cref{lemma:asw_etale_torsor}. The assignment
$[\vec g]\mapsto[T_{\vec g}]$ is a canonical isomorphism of abelian groups
\[
W_r(M)/\wp W_r(M)\;\xrightarrow{\ \cong\ }\;
H^1_{\mathrm{\acute et}}
\bigl(\Spec M,(\Z/p^r\Z)_{\Spec M}\bigr).
\]
Thus this quotient classifies $\Z/p^r\Z$-torsors over $\Spec M$.

If $[\vec g]$ has order $p^s$, where $0\leq s\leq r$, then $T_{\vec g}$
has $p^{r-s}$ connected components, each noncanonically isomorphic to
$\Spec L$ for the same cyclic extension $L/M$ of degree $p^s$.
Let $H\subseteq\Z/p^r\Z$ be the unique subgroup of order $p^s$.
Each component is an $H$-torsor, and $T_{\vec g}$ is induced from any one
of them. In particular, $T_{\vec g}$ is connected if and only if
$[\vec g]$ has order $p^r$. Every cyclic extension of $M$ of degree
dividing $p^r$ occurs in this way.
\end{theorem}

\begin{proof}
Let $M^s$ be a separable closure of $M$, and put $\Gamma=\Gal(M^s/M)$.
The triangular equations in the proof of
\Cref{lemma:asw_etale_torsor} show that $\wp$ is surjective on $W_r(M^s)$
and has kernel $W_r(\F_p)\cong\Z/p^r\Z$. Choose $\vec u\in W_r(M^s)$
with $\wp\vec u=\vec g$. Under the torsor--character correspondence
of \Cref{cor:abelian_equivilance_fundamental_torsors}, $T_{\vec g}$
corresponds to
\[
\chi_{\vec g}\colon\Gamma\longrightarrow W_r(\F_p),
\qquad
\chi_{\vec g}(\sigma)=\sigma(\vec u)\ominus\vec u.
\]
This is the character of the connecting class $[T_{\vec g}]$ in
\Cref{eq:common_torsor_connecting_map}.

We first check that this map is well defined. Any other solution of
$\wp\vec u=\vec g$ is $\vec u\oplus\vec c$ for some
$\vec c\in W_r(\F_p)$. Since $\Gamma$ fixes $\vec c$, this change does not
alter $\chi_{\vec g}$. Similarly, replacing $\vec g$ by
$\vec g\oplus\wp\vec v$, with $\vec v\in W_r(M)$, replaces $\vec u$ by
$\vec u\oplus\vec v$ and again leaves the character unchanged. If
$\chi_{\vec g}=0$, then $\vec u$ is fixed by $\Gamma$, so
$\vec u\in W_r(M)$ by \Cref{eq:witt_invariants_componentwise}. Hence
$[\vec g]=0$, and the map is injective.

It remains to prove surjectivity. We claim that
\[
H^1_{\mathrm{cont}}\bigl(\Gamma,W_r(M^s)\bigr)=0.
\]
For $r=1$, this is additive Hilbert~90. For $r>1$,
\Cref{lemma:witt_bookkeeping}(i) gives a short exact sequence of
$\Gamma$-modules
\[
0\longrightarrow M^s\xrightarrow{V^{r-1}}W_r(M^s)
\xrightarrow{t}W_{r-1}(M^s)\longrightarrow0.
\]
The associated long exact sequence contains
\[
H^1_{\mathrm{cont}}(\Gamma,M^s)
\longrightarrow
H^1_{\mathrm{cont}}\bigl(\Gamma,W_r(M^s)\bigr)
\longrightarrow
H^1_{\mathrm{cont}}\bigl(\Gamma,W_{r-1}(M^s)\bigr).
\]
The outer groups vanish by additive Hilbert~90 and induction, so the middle
group vanishes. The long exact sequence associated with
\Cref{eq:asw_lang_sequence} now shows that the connecting map is surjective.
This proves the claimed isomorphism.

Now suppose $[\vec g]$ has order $p^s$, and put
$H=\operatorname{im}(\chi_{\vec g})$. The isomorphism preserves orders,
so $\chi_{\vec g}$ has order $p^s$. Since $H$ is cyclic, the order of the
character equals $|H|$. By \Cref{theorem:equivalence_fundamental_covers},
the connected components of $T_{\vec g}$ correspond to the $\Gamma$-orbits
on its geometric fiber. Choosing $\vec u$ identifies that fiber with
$W_r(\F_p)$, with $\Gamma$ acting by translation through $\chi_{\vec g}$.
The orbits are the cosets of $H$, so there are $p^{r-s}$ components.
Each geometric point has stabilizer $\ker(\chi_{\vec g})$ in $\Gamma$.
Thus each component is isomorphic to $\Spec L$, where
$L=(M^s)^{\ker(\chi_{\vec g})}$ is cyclic of degree $|H|=p^s$ over $M$.
The subgroup $H$ preserves each component and acts simply transitively on
its geometric fiber. Each component is therefore an $H$-torsor, and
inducing along $H\hookrightarrow W_r(\F_p)$ recovers $T_{\vec g}$.

Conversely, let $L/M$ be cyclic of degree $p^s$ dividing $p^r$, and embed
$L$ in $M^s$. Choose an injection $\Gal(L/M)\hookrightarrow W_r(\F_p)$
and compose it with the
restriction map $\Gamma\to\Gal(L/M)$. By the surjectivity proved above, the
resulting character is $\chi_{\vec g}$ for some $\vec g\in W_r(M)$. The
fixed field of its kernel is $L$, so the connected components of $T_{\vec g}$ are
isomorphic to $\Spec L$.
\end{proof}

% [PROPOSED REWRITE, for comparison with the theorem and proof above]
{\color{green!50!black}
\begin{theorem*}[Proposed rewrite of \Cref{theorem:asw_theory}]
Let $M$ be a field of characteristic $p$, and let $r\geq1$. For
$\vec g\in W_r(M)$ put
$T_{\vec g}:=\Spec M[\vec X]/(\wp\vec X\ominus\vec g)$.
\begin{enumerate}[label=(\roman*)]
  \item The map $[\vec g]\mapsto[T_{\vec g}]$ is a group isomorphism
  \[
  W_r(M)/\wp W_r(M)\xrightarrow{\ \cong\ }
  H^1_{\mathrm{\acute et}}\bigl(\Spec M,(\Z/p^r\Z)_{\Spec M}\bigr).
  \]
  The target is \'etale cohomology with coefficients in the constant group
  scheme $(\Z/p^r\Z)_{\Spec M}\cong W_r(\F_p)_{\Spec M}$
  (\Cref{example:witt_of_Fp}). Its elements are the classes of
  $\Z/p^r\Z$-torsors over $\Spec M$ (see the note in
  \Cref{section:torsors_and_cohomology_note}).
  \item Suppose $[\vec g]$ has order $p^s$, and let $H\subseteq\Z/p^r\Z$
  be the unique subgroup of order $p^s$. Then $T_{\vec g}$ is a disjoint
  union of $p^{r-s}$ copies of $\Spec L$, where $L/M$ is cyclic of degree
  $p^s$. Each copy is an $H$-torsor, and $T_{\vec g}$ is induced from it.
  In particular, $T_{\vec g}$ is connected if and only if $[\vec g]$ has
  order $p^r$.
  \item Every cyclic extension of $M$ of degree dividing $p^r$ is such
  an $L$.
\end{enumerate}
\end{theorem*}

\begin{proof}[Proposed proof]
\emph{(i), injectivity.} Apply \Cref{lemma:asw_etale_torsor} to the
universal vector over $\F_p[g_0,\dots,g_{r-1}]$: the map
$\wp:\bW_r\to\bW_r$ is finite \'etale and surjective, with kernel
$W_r(\F_p)\cong\Z/p^r\Z$. So \Cref{eq:common_torsor_connecting_map}
gives an injective homomorphism $[\vec g]\mapsto[T_{\vec g}]$.

\emph{(i), surjectivity.} The long exact sequence in
\Cref{appendix:common_group_scheme_construction} continues with
$H^1(M,\Z/p^r\Z)\to H^1_{\mathrm{\acute et}}(\Spec M,\bW_r)$. So it is
enough to show that $H^1_{\mathrm{\acute et}}(\Spec M,\bW_r)=0$. For $r=1$
this is additive Hilbert~90 \cite[X, \S1]{serreLF}. For $r>1$, the sequence
$0\to\bG_a\xrightarrow{V^{r-1}}\bW_r\xrightarrow{t}\bW_{r-1}\to0$
(\Cref{prop:witt_multiplication_by_p}) is exact. Its outer terms have zero
$H^1$ by induction, so the middle one does too.

\emph{(ii).} Let $\chi:\Gal(M^s/M)\to\Z/p^r\Z$ be the character of
$T_{\vec g}$ (\Cref{cor:abelian_equivilance_fundamental_torsors}). The map
in (i) is a group isomorphism, so $\chi$ has order $p^s$. A subgroup of
$\Z/p^r\Z$ is cyclic, so the order of $\chi$ equals the size of its image.
Hence $\operatorname{im}\chi=H$. By
\Cref{theorem:equivalence_fundamental_covers}, $T_{\vec g}$ is induced from
the connected $H$-torsor $\Spec L$, where $L$ is the fixed field of
$\ker\chi$ and $[L:M]=|H|=p^s$. An induced torsor has
$[\Z/p^r\Z:H]=p^{r-s}$ components.

\emph{(iii).} Let $L/M$ be cyclic of degree $p^s\mid p^r$. Choose an
injection $\Gal(L/M)\hookrightarrow\Z/p^r\Z$. The composite
$\Gal(M^s/M)\to\Z/p^r\Z$ is a character. By (i) it comes from some
$\vec g$, and by (ii) the components of $T_{\vec g}$ are $\Spec L$.
\end{proof}
}

\begin{cor}\label{cor:asw_class_order_and_first_component}
Let $M\supseteq\F_p$ and let $\vec g\in W_r(M)$ have $g_0=0$. Then
\[
p^{r-1}[\vec g]=0 \qquad\text{in } W_r(M)/\wp W_r(M).
\]
Consequently, a class of
order $p^r$ has nonzero zeroth component in every representative.
\end{cor}

\begin{proof}
By \Cref{eq:witt_p_to_r_minus_1},
$p^{r-1}\vec g=(0,\dots,0,g_0^{p^{r-1}})=0$ already in $W_r(M)$.
\end{proof}

\subsection{Reduced polynomial representatives}

\begin{definition}\label{definition:reduced_polynomial_vector}
A vector $\vec g=(g_0,\dots,g_{r-1})\in W_r(k((x)))$ is a
\textbf{reduced polynomial vector} if it is reduced and every
$g_j\in x^{-1}k[x^{-1}]$.
\end{definition}

\begin{lemma}[Surjectivity on integral Witt vectors]
\label{lemma:wp_surjective_on_integral_witt}
If $k$ is algebraically closed, then
\[
\wp:W_r(k[[x]])\longrightarrow W_r(k[[x]])
\]
is surjective.
\end{lemma}

\begin{proof}
Proceed by induction on $r$. For $r=1$, given
$h=\sum_{i\geq0}b_ix^i$, choose $u_0$ with
$u_0^p-u_0=b_0$ and define the remaining coefficients recursively by
$u_i=u_{i/p}^p-b_i$, with $u_{i/p}=0$ when $p\nmid i$. Then
$u^p-u=h$.

For $r>1$, choose by induction $\vec u'\in W_{r-1}(k[[x]])$ with
$\wp\vec u'=t(\vec h)$ and lift it to
$\widetilde u=(\vec u',0)$. By \Cref{lemma:witt_bookkeeping}(i),
$\wp\widetilde u\ominus\vec h=V^{r-1}(g)$ for some $g\in k[[x]]$.
Choose $w\in k[[x]]$ with $\wp w=-g$. Then
\Cref{lemma:witt_bookkeeping}(ii) gives
$\wp(\widetilde u\oplus V^{r-1}w)=\vec h$.
\end{proof}

\begin{prop}[Reduced polynomial representatives]
\label{prop:reduced_polynomial_representatives}
Every class in $W_r(k((x)))/\wp W_r(k((x)))$ has a reduced polynomial
representative.
\end{prop}

\begin{proof}
Proceed by induction on $r$. For $r=1$, write a Laurent series as the sum of
its principal part and an element of $k[[x]]$. The latter can be removed by
\Cref{lemma:wp_surjective_on_integral_witt}. If the remaining principal part
has leading term $cx^{-m}$ with $p\mid m$, subtract
\[
\wp(c^{1/p}x^{-m/p})=cx^{-m}-c^{1/p}x^{-m/p}.
\]
This strictly decreases the pole order; iteration produces zero or a pole
order prime to $p$.

For the inductive step, first reduce the truncation of the given vector and
lift the change of representative to length $r$. The first $r-1$ components
are then reduced polynomials. By
\Cref{lemma:witt_bookkeeping}(iii), subtracting an element of the form
$\wp(V^{r-1}z)=V^{r-1}(\wp z)$ changes only the last component. Apply the
length-one procedure to that component: remove its integral part and then
successively remove leading terms whose pole orders are divisible by $p$.
This leaves all earlier components unchanged and produces a reduced polynomial
vector representing the original class.
\end{proof}
